%%%%%%%%%%%%%%%%%%%%%%%%%%%%%%%%%%%%%%%
% Cover Letter - AKQA, Product Design Director, Berlin
%%%%%%%%%%%%%%%%%%%%%%%%%%%%%%%%%%%%%%%

\documentclass[]{cover-fira}
\usepackage{fancyhdr}

\pagestyle{fancy}
\fancyhf{}

\rfoot{Page \thepage \hspace{0pt}}
\thispagestyle{empty}
\renewcommand{\headrulewidth}{0pt}
\begin{document}

%%%%%%%%%%%%%%%%%%%%%%%%%%%%%%%%%%%%%%
%     TITLE NAME
%%%%%%%%%%%%%%%%%%%%%%%%%%%%%%%%%%%%%%
\namesection{}{\Huge{Marko Rosić}}{%
  \href{mailto:marko@rosic.net}{marko@rosic.net} | +381 60 585 44 55 |
  \urlstyle{same}\href{https://linkedin.com/in/markorosic}{LinkedIn}
}

%%%%%%%%%%%%%%%%%%%%%%%%%%%%%%%%%%%%%%
%     MAIN COVER LETTER CONTENT
%%%%%%%%%%%%%%%%%%%%%%%%%%%%%%%%%%%%%%

\currentdate{\today}
\lettercontent{Dear AKQA hiring team,}

\lettercontent{I am applying for the Product Design Director role in Berlin. What draws me to AKQA is not the client roster — it is the specific challenge of making brand intent implementable at scale. The hardest design systems problems I have worked on are not about visual language; they are about defining what can be reused, what must vary per context, and how to make those decisions legible to engineers across multiple teams. That is as true for a luxury brand's ecosystem as it is for a multi-market SaaS product.}

\lettercontent{My background is in-house product, which is honest. I have not worked agency-side. What I bring from that background:}

{\raggedright\fontspec{Fira Sans}\fontsize{11pt}{13pt}\selectfont
\begin{itemize}
    \item \textbf{Design systems at multi-brand scale:} I have built two large token-based component libraries from scratch — one that reduced development time by 30\% across 10+ websites at Catena Media, and one that cut time-to-market by 25\% across a multi-brand portfolio at Gentoo Media. Component ownership, cross-brand consistency, and engineering handoff are where I spend the most time and have the most concrete results.
    \item \textbf{Director-level team leadership:} at Gentoo Media I built and led a team of seven, including two team leads who each ran their own sub-team. I set design direction, own performance, and build the structure so the team can work without me in the room.
    \item \textbf{AI-driven design workflow:} I use Claude Code to build working prototypes and Figma plugins. This is not a workflow experiment — it is how I deliver. It means I stay close to what is technically achievable and move faster from concept to validated design.
    \item \textbf{Evidence-based practice:} I bring analytics, A/B testing, and user research into every design process. I do not design for brand aesthetics alone; I design for user behaviour and measure the outcome.
\end{itemize}\par}
\vspace{6pt}

\lettercontent{I am aware the Director role at AKQA involves client management and executive presentation — skills I have developed in an in-house context but have not exercised in an agency setting. I am applying because AKQA Berlin's technology client work is where I see the tightest overlap with the kind of design infrastructure problems I have built my practice around.}

\lettercontent{I am based in Belgrade (CET) and open to relocation to Berlin. I would like to learn more about the role and the Berlin studio's current focus.}

\begin{flushright}
\closing{Kind regards,}

\signature{Marko Rosić}
\end{flushright}
\end{document}
